\chapter{Ausgangslage}
Schachcomputer sind faszinierend: Mit erstaunlicher Präzision berechnen sie Varianten, bewerten Stellungen und finden
Züge, die für Menschen manchmal schwer nachvollziehbar sind. Durch mathematische Modelle und effiziente Suchverfahren haben
sie das Spiel in den letzten Jahrzehnten so stark \enquote{perfektioniert}, dass sie heute selbst die besten menschlichen
Spieler übertreffen.

Im modernen Schach sind Schachcomputer daher nicht mehr wegzudenken. Sie werden zur Vorbereitung, zum Analysieren von
Partien und zur Verbesserung des eigenen Spiels eingesetzt.

Als Schachliebhaber möchten wir, Dario und Romeo, uns mit diesem Thema beschäftigen. Unser Ziel ist es zu verstehen,
wie ein Schachcomputer aufgebaut ist, wie er vorgeht und anschliessend einen einfachen Schachcomputer selbst zu
entwickeln.

Uns ist bewusst, dass unser Projekt nicht mit kommerziellen Programmen oder Systemen konkurrieren kann, hinter denen
ganze Teams und grosse Communitys stehen. Dennoch finden wir es spannend, wie weit man mit einfachen Mitteln kommen kann.
